% Statement-only: night:P1.1
\begin{theorem}\label{night:P1.1}
\StatementRecord{night:P1.1}{proved}{}{Definitions and the displayed constructions}
The finite homology monodromies have orders $5,4$, respectively, and
characteristic polynomials
$(T-1)^2\Phi _5(T)$ and $(T-1)^2(T+1)^2(T^2+1)$.
For $N=M_0-I$ we have $N^2=0$, $\operatorname{rank}N=3$, and
nonzero Smith factors $1,1,1$.  Moreover
\[
 \Lambda^{\langle M_1,M_2\rangle}=0,\qquad
 \Lambda_{\langle M_1,M_2\rangle}=\mathbb Z,
\]
where the quotient homomorphism is $\ell(x)=x_0$.
On the cohomology lattice the corresponding groups are $\mathbb Ze_0$
and $\mathbb Z/2$.  Every nonidentity power of each specified affine
action $\tau_i$ is free.
\end{theorem}

% Statement-only: night:P1.4
\begin{theorem}\label{night:P1.4}
\StatementRecord{night:P1.4}{proved}{}{Definitions and the displayed constructions}
The genus-two pencil has one nodal fibre away from $0,\infty$, at
$\lambda_*=-256/3125$.  Good reduction at $0,\infty$ requires root
orders $8,5$.  On this root orbifold its Hodge bundle is
$\mathcal H=T^3\oplus T$, where
\[
T=\mathcal O\bigl((2/5)[\infty]-(3/8)[0]\bigr),\qquad
\deg T=1/40,
\]
and $T^3$ is the line generated meromorphically by $dx/y$.
\end{theorem}

% Statement-only: night:P1.5
\begin{theorem}\label{night:P1.5}
\StatementRecord{night:P1.5}{proved}{}{Definitions and the displayed constructions}
For the Hodge line $\ell_0\subset\mathcal H$ of the genus-two pencil,
the image of $\ell_0\otimes\mathcal H\to\bigwedge^2\mathcal H$
has rank one. The wedge map from
$\ell_0\otimes(\mathcal H_{\mathrm{dR}}/\ell_0)$, followed by
projection modulo the nondegenerate polarization bivector, is
injective and has rank-three image. This image defines the bundle
$P$ used in the period construction.
\end{theorem}

% Statement-only: night:P1.7
\begin{theorem}\label{night:P1.7}
\StatementRecord{night:P1.7}{proved}{}{\cref{maps:src:F.3.5}, \cref{night:P1.2}}
The integral $H^1$ monodromy of the pencil has a companion marking with
finite pair $(C_f^{-1},C_g)$, and cusp product $C_fC_g^{-1}$.
Its exterior-square quotient is consequently the integral quotient
representation appearing in P1.2.
\end{theorem}

% Statement-only: night:P1.8
\begin{theorem}\label{night:P1.8}
\StatementRecord{night:P1.8}{proved}{}{\cref{night:P1.5}}
With the corrected rank-three bundle $P$ of P1.5, the supplied rank-six
system admits a holomorphic lifted family
$F_c=(\sigma+ic\alpha e)(P)$ on the punctured base.
For a specified splitting $\sigma$ there is a real analytic function
$h$ with finite attained minimum $h_{\min}$, and the exact domain of
real parameters giving $F_c\cap\overline{F_c}=0$ on the completed
base, including both finite root points and excluding only the cusp
before its toroidal filling, is $c>-h_{\min}$.
At the cusp its periods have the form
$Z(q)=(\log q/(2\pi i))I_3+C(q)$ with $C$ holomorphic and bounded.
The construction extends to the two finite root charts and descends to
orders $5,4$ on the rank-six system.
\end{theorem}

% Statement-only: night:P1.15
\begin{theorem}\label{night:P1.15}
\StatementRecord{night:P1.15}{proved}{}{\cref{night:P1.1}, \cref{night:P1.8}}
The rank-six flat extension, as a complex flat bundle with its invariant
line, has the following explicit algebraic connection model.  Set
$\omega_k=x^kdx/y$, $0\leq k\leq3$, and use the quotient frame
\[
(w_2,w_1,w_0,w_{-1},w_{-2})
 =(\omega_{01},\omega_{02},\omega_{03},\omega_{13},\omega_{23}),
\qquad \omega_{12}=-3\omega_{03}.
\]
These are meromorphic frames of
$(L^2,L,\mathcal O,L^{-1},L^{-2})$, where $L=T^2$ on the root
orbifold of orders $5,4$.  If a frame is written as a row of vectors,
its derivative is right multiplication by
\[
\Gamma_6=\begin{pmatrix}0&b\\0&\Gamma_W\end{pmatrix},\qquad
b=\frac1{\lambda-\lambda_*}
 \left(-\frac{25}{6\lambda^2},-\frac{125}{48\lambda},0,1,0\right),
\tag{1.15.1}
\]
where
\[
\resizebox{\textwidth}{!}{$\displaystyle
\Gamma_W=\frac1{3125\lambda+256}
\begin{pmatrix}
-8(625\lambda+48)/\lambda&-8&10\lambda&75\lambda^2/2&0\\
-125/(2\lambda)&-4(625\lambda+48)/\lambda&8&0&75\lambda^2/2\\
300/\lambda^2&-375/\lambda&0&-48&60\lambda\\
120/\lambda^3&0&125/(2\lambda)&4(625\lambda+48)/\lambda&-8\\
0&120/\lambda^3&50/\lambda^2&-125/(2\lambda)&8(625\lambda+48)/\lambda
\end{pmatrix}.$}
\tag{1.15.2}
\]
This model specifies a canonical algebraic splitting, denoted
$\sigma_b$.  The compatible contraction used for the real parameter is
$\alpha=(625\pi i/128)\alpha_{\rm alg}$, where
$\alpha_{\rm alg}(w_0)=4/3$ and it vanishes on $w_1,w_2$.
\end{theorem}

% Statement-only: night:P1.18
\begin{theorem}\label{night:P1.18}
\StatementRecord{night:P1.18}{proved}{}{\cref{night:P1.15}; the rational connection and root-gauge definitions displayed in \cref{night:P1.16}; \cref{night:P1.8}; \cite[Lemma~2.7, pp.~42--43]{TeschlODE2012}; \cite[\S1.9(iii), (1.9.30)--(1.9.31)]{DLMFCauchy2026}}
In the algebraic splitting $\sigma_b$, the function $h_b$ has its
unique global minimum at the finite-root point $\lambda=\infty$.
The minimum is the exact Frobenius period constant
\[
 h_{b,\min}=-\frac{96}{625\pi}\operatorname{Re}(XH_\infty)_{03},
 \qquad
 h_{b,\min}=-1.0945402503114027801162682630549457037\ldots,
\tag{1.18.1}
\]
where $X$ is the normalized integral intertwiner of P1.15 and
$H_\infty$ is the regular infinity-root limit of its fundamental
solution, normalized to be the identity at $\lambda=1$ before
multiplication by $X$.  Formula (1.18.1), the rational connection,
and those two normalizations specify the constant exactly.
A fresh certified enclosure has radius less than $4\cdot10^{-76}$
about the full stored midpoint in
\path{code/paper_v/output/fresh_frobenius.json}.
Consequently the exact domain of completed transverse periods is
\[
 c>-h_{b,\min},\qquad
 \operatorname{Im}\xi>-h_{b,\min}
\]
for the real and complex parameters, respectively.
The member $c=10$ used in the numerical period checks is therefore
admissible on the whole completed base.

For the point-normalized splitting of P1.8, continued from
$1$ to $\lambda_*/2$ by the polygon
$1,1+i,\lambda_*/2+i,\lambda_*/2$, the exact comparison constant
$a_{\rm split}$ of (1.15.4) has certified value
\[
\begin{split}
 \operatorname{Re}a_{\rm split}
   &=-0.7099885679126427771723473455012308149\ldots,\\
 \operatorname{Im}a_{\rm split}
   &=-0.9134017059364153362143229223828373269\ldots.
\end{split}
\]
Thus the minimum in that splitting is
$-0.1811385443749874439019453406721083768\ldots$,
attained at the same base point.  Its real parameter interval is
\[
 c>0.1811385443749874439019453406721083768\ldots.
\]
The full parameter comparison remains
$\xi_b=\xi_{\rm original}-a_{\rm split}$; both its real and imaginary
parts are necessary to identify the two normalized families.
\end{theorem}
